\chapter{Background \& Related Work}
\label{chapter:background}


\section{Traditional Simultaneous Interpretation Systems}

Simultaneous interpretation has long been a cornerstone of multilingual communication in international conferences, governmental institutions, and large-scale events. In traditional settings, interpretation is performed by highly trained human interpreters who listen to a speaker in a source language and deliver a near-synchronous spoken translation in a target language. This process requires a high level of linguistic expertise, cognitive load management, and sustained concentration, making professional interpreters a scarce and costly resource \cite{gencheva2023features}.

In addition to human expertise, traditional simultaneous interpretation relies on a specialized physical infrastructure designed to ensure low-latency audio transmission and clear separation between source and target language channels. Interpreters typically work inside sound-isolated booths, where they receive the original speech through headphones and produce the translated output via microphones. These booths are engineered to minimize external noise and prevent audio feedback, thereby maintaining translation accuracy and listener comfort.

The interpreter audio is routed through dedicated interpretation consoles that allow control over input channels, volume levels, and language selection. From there, the translated signal is transmitted using wired or wireless distribution systems to receivers used by audience members. Listeners select their preferred language channel on handheld receivers or integrated conference devices, enabling simultaneous access to multiple language streams \cite{conference_interpreting_new_tech}. The components of such a traditional interpretation setup are visualized in Figure~\ref{fig:traditional_interpretation_system}.

\begin{figure}[!htbp]
  \centering
  \includegraphics[width=\textwidth]{img/figures/traditional_interpretation_system.drawio.png}
  % \footnotemark creates the red number in the caption
  \caption[Traditional Simultaneous Interpretation System Components]{Traditional simultaneous interpretation system components. Adapted from \footnotemark}
  \label{fig:traditional_interpretation_system}
\end{figure}

% \footnotetext places the link at the bottom of the current page
\footnotetext{\url{https://media.boschsecurity.com/fs//media/pt/pb/media/products_1/conference_systems_1/integrus_brochure.pdf}}

This tightly integrated hardware ecosystem ensures reliable, low-latency delivery of interpreted speech but comes with significant logistical requirements \cite{ozkaya2023transformative}. Deploying a traditional interpretation system involves transporting and installing booths, consoles, transmitters, receivers, and cabling, often requiring specialized technicians. Furthermore, each additional target language increases system complexity, as it necessitates additional interpreters, booths, and audio channels. As a result, scalability is limited, particularly for events with many languages or dynamically changing interpretation needs.

From an economic perspective, traditional simultaneous interpretation incurs substantial recurring costs. These include interpreter fees, equipment rental, installation labor, and on-site technical support. Such costs can be prohibitive for smaller organizations or informal settings, restricting access to high-quality multilingual communication. Moreover, reliance on physical infrastructure constrains deployment flexibility, making rapid or remote interpretation scenarios difficult to support \cite{rayaa2022remoteSI}.

Despite these limitations, traditional interpretation systems remain the industry standard due to their proven reliability and high output quality. They therefore provide an important baseline against which alternative approaches must be evaluated. In particular, any AI-based interpretation system aiming for real-world applicability must account for similar requirements in terms of latency, audio routing, and hardware compatibility \cite{dhawan2022speech2speechchallenges}. The feasibility of integrating AI-driven interpretation pipelines with professional audio equipment is therefore a critical consideration and is revisited later in this thesis in Section~\ref{chapter:hardware}.


\section{Core Technologies for Real-Time Speech Interpretation}

This section provides background on the core technologies that underpin modern real-time speech interpretation systems. In particular, it reviews the roles and operational characteristics of Automatic Speech Recognition, Neural Machine Translation, and Text-to-Speech synthesis when deployed under strict real-time constraints. Rather than focusing on model architectures or training methodologies, the discussion emphasizes streaming behavior, latency considerations, and system-level challenges that arise when these components are combined into an end-to-end interpretation pipeline. This technological foundation establishes the context for the system design choices presented in Chapter~\ref{chapter:methodology}.

\subsection{Automatic Speech Recognition in Streaming Settings}

Automatic Speech Recognition (ASR) refers to the process of converting spoken language into written text using computational systems \cite{nayeem2025asr_survey}. In the context of multilingual communication, ASR serves as the primary interface between continuous audio input and downstream language processing components \cite{8252104}. For real-time speech interpretation, ASR systems must operate under strict temporal constraints, producing textual output incrementally as speech unfolds rather than after complete utterances are available. A comparison between batch-based and streaming automatic speech recognition, highlighting incremental processing and partial hypothesis generation in streaming settings is visualised in Figure~\ref{fig:streaming_vs_batch_asr}.

\begin{figure}[!htbp]
  \centering
  \includesvg[width=\textwidth]{img/figures/streaming_vs_batch_asr.drawio.svg}
  \caption[Batch-based and Streaming Automatic Speech Recognition Comparison]{Batch-based and Streaming Automatic Speech Recognition Comparison}
  \label{fig:streaming_vs_batch_asr}
\end{figure}

Traditional ASR systems have often been designed for batch or offline processing, where an entire audio recording is first captured and then transcribed in a single pass. While such approaches can leverage full acoustic context to improve transcription accuracy, they introduce delays that are incompatible with simultaneous interpretation scenarios. In contrast, streaming ASR systems process audio input continuously and emit recognition results in near real time. This distinction is fundamental for applications that require immediate downstream processing, such as live translation or captioning.

A defining characteristic of streaming ASR is the generation of partial hypotheses. As audio is processed incrementally, the recognizer produces provisional text outputs that may be revised as additional acoustic context becomes available. These partial hypotheses allow downstream components to react early to spoken content, thereby reducing overall latency. However, they also introduce instability, as previously emitted text segments may change or be retracted. Streaming ASR systems therefore balance two competing objectives: minimizing latency while maintaining sufficient stability in the generated transcripts \cite{weller2021streamingST}.

In real-time interpretation settings, this trade-off between responsiveness and stability becomes particularly critical. Partial hypotheses that are forwarded too aggressively may lead to semantic inconsistencies or repeated processing in downstream stages. Conversely, delaying output until hypotheses are fully stable increases end-to-end latency and undermines the simultaneity of interpretation \cite{grissom2014realtime}. As a result, streaming ASR is not merely a faster variant of offline recognition but a qualitatively different operating mode with distinct design considerations.

Another important distinction between streaming and batch ASR lies in how speech boundaries are handled. Offline systems can rely on explicit sentence boundaries or complete utterances, whereas streaming systems must operate without prior knowledge of sentence length or structure. Instead, they infer boundaries implicitly from acoustic cues, pauses, or internal confidence measures \cite{weller2021streamingST}. This behavior directly affects how recognized speech can be segmented and forwarded for translation in real time.

From a system perspective, streaming ASR is typically integrated as a continuous service that accepts audio chunks of fixed duration and returns updated recognition hypotheses asynchronously. This incremental interaction model enables low-latency processing but requires careful coordination with downstream components that consume ASR output. In particular, downstream stages must be designed to tolerate hypothesis revisions and to distinguish between provisional and finalized recognition results.


\subsection{Neural Machine Translation for Spoken Language}

Neural Machine Translation (NMT) refers to the use of neural network–based models to automatically translate text from a source language into a target language \cite{stahlberg2020nmt_review}. In contrast to earlier rule-based or statistical approaches, modern NMT systems model translation as a sequence-to-sequence transformation and have become the dominant paradigm in both academic research and commercial deployment. Within real-time speech interpretation systems, NMT serves as the core linguistic transformation stage that converts recognized speech content into translated text suitable for immediate delivery or speech synthesis \cite{sperber2020endtoendpromise}.

Most established NMT systems are designed for text-based translation scenarios, such as documents, subtitles, or conversational messages, where complete sentences or paragraphs are available prior to translation. In such settings, the translation model can leverage full contextual information, including sentence boundaries and long-range dependencies, to maximize output quality. However, real-time interpretation imposes fundamentally different constraints. Instead of well-formed textual input, NMT systems receive short, incrementally generated speech segments that may represent incomplete clauses or evolving utterances \cite{chen2024metamorphosisMT}.

This distinction between document-level text translation and speech segment translation has important implications. In spoken language interpretation, translation must often be performed on partial and temporally constrained input, without access to the full sentence context. As a result, translation decisions are made under uncertainty, and the system must balance early output generation against the risk of semantic incompleteness. These constraints differentiate real-time NMT from traditional text translation tasks and require system-level strategies that prioritize responsiveness alongside linguistic coherence.

A central challenge in spoken-language NMT is the trade-off between latency and translation quality. Translating shorter segments reduces delay and improves simultaneity but may degrade fluency or lead to awkward phrasing due to missing context. Conversely, aggregating longer segments can improve translation accuracy but increases end-to-end latency, which is undesirable in live interpretation scenarios. This trade-off is inherent to real-time systems and cannot be resolved solely through model improvements; instead, it must be managed through careful segmentation and pipeline design.

Segment-based translation introduces additional challenges related to consistency and continuity. Since continuous segments are translated instantly in a connected stream, maintaining coherent terminology, pronoun references, and discourse flow across segments becomes more difficult than in document-level translation \cite{dhawan2022speech2speechchallenges}. Furthermore, variations in segment length and timing can lead to uneven translation pacing, which must be handled gracefully by downstream components such as text-to-speech synthesis. These challenges highlight the importance of viewing NMT as one component within a broader real-time system rather than as an isolated module.

Despite these limitations, segment-level NMT remains the most practical approach for real-time interpretation due to its compatibility with streaming pipelines and asynchronous execution models. By translating speech incrementally, systems can deliver translated content with minimal delay while preserving acceptable semantic quality. As summarized in Table~\ref{tab:nmt_text_vs_speech}, the operational requirements of spoken-language translation differ substantially from those of traditional text translation, motivating specialized system-level adaptations for real-time use cases.


\begin{table}[!htbp]
  \centering
  \caption[Text translation and segment-based translation comparison]{Comparison between document-level text translation and segment-based translation for spoken language}
  \label{tab:nmt_text_vs_speech}
  \begin{adjustbox}{max width=\textwidth, keepaspectratio=false}
    \begin{tabular}{p{4.3cm} p{7.5cm} p{7.5cm}}
      \toprule
      \textbf{Aspect} & \textbf{Text Translation} & \textbf{Spoken-Language Translation} \\
      \midrule
      Input characteristics 
      &
      \begin{itemize}[nosep,leftmargin=*]
        \item Complete sentences or full documents
        \item Static input available before translation
        \item Explicit sentence boundaries
      \end{itemize}
      &
      \begin{itemize}[nosep,leftmargin=*]
        \item Incremental speech segments
        \item Input evolves over time
        \item Implicit and dynamically inferred boundaries
      \end{itemize} \\
      \midrule
      Context availability 
      &
      \begin{itemize}[nosep,leftmargin=*]
        \item Global context across sentences or paragraphs
        \item Long-range dependencies accessible
        \item Stable discourse structure
      \end{itemize}
      &
      \begin{itemize}[nosep,leftmargin=*]
        \item Limited local context
        \item Partial or incomplete semantic information
        \item No access to future utterances
      \end{itemize} \\
      \midrule
      Latency requirements 
      &
      \begin{itemize}[nosep,leftmargin=*]
        \item Latency typically non-critical
        \item Offline or asynchronous processing acceptable
        \item Quality prioritized over speed
      \end{itemize}
      &
      \begin{itemize}[nosep,leftmargin=*]
        \item Strict real-time constraints
        \item Low end-to-end delay required
        \item Latency directly affects usability
      \end{itemize} \\
      \midrule
      Translation granularity 
      &
      \begin{itemize}[nosep,leftmargin=*]
        \item Sentence-level or paragraph-level units
        \item Coherent, complete linguistic structures
      \end{itemize}
      &
      \begin{itemize}[nosep,leftmargin=*]
        \item Clause-level or short segment-level units
        \item Semantically incomplete fragments possible
      \end{itemize} \\
      \midrule
      Stability and revision risk 
      &
      \begin{itemize}[nosep,leftmargin=*]
        \item No revision after translation
        \item Input text remains unchanged
      \end{itemize}
      &
      \begin{itemize}[nosep,leftmargin=*]
        \item Risk of revision due to evolving speech
        \item Dependence on ASR hypothesis stability
      \end{itemize} \\
      \bottomrule
    \end{tabular}
  \end{adjustbox}
\end{table}

\vspace{3cm}
\subsection{Text-to-Speech in Live Interpretation Scenarios}

Text-to-Speech (TTS) systems are responsible for converting textual representations of language into intelligible spoken audio \cite{tan2021tts_survey}. In the context of real-time speech interpretation, TTS constitutes the final stage of the processing pipeline, directly affecting the listener’s perception of responsiveness, fluency, and overall system quality. Unlike offline speech synthesis applications, live interpretation imposes strict constraints on latency and audio delivery, requiring TTS systems to operate incrementally and with minimal delay.

In traditional applications such as audiobook narration or voice assistants, speech synthesis can be performed on complete sentences or paragraphs, allowing models to optimize for prosody, expressiveness, and naturalness. However, in simultaneous interpretation, TTS systems receive short, segment-level translation outputs that may represent incomplete clauses or evolving utterances. As a result, synthesis must begin promptly upon receiving each segment, even when full linguistic context is unavailable \cite{popescu-belis2025s2streview}. This operational mode fundamentally differentiates live interpretation TTS from conventional speech synthesis tasks.

Historically, TTS systems were based on concatenative or parametric approaches, which relied on pre-recorded speech units or handcrafted acoustic models. While these methods offered predictable latency and low computational overhead, they often produced unnatural or monotonous speech \cite{cohn2020perceptionTTS}. Recent advances in neural TTS have significantly improved naturalness and speaker expressiveness, enabling synthetic speech that closely resembles human voices. These improvements have facilitated the adoption of neural TTS in real-time applications, including interpretation systems, despite their higher computational complexity \cite{cohn2020perceptionTTS}. To highlight the differences between traditional and neural text-to-speech approaches in the context of live interpretation, Table~\ref{tab:tts_traditional_vs_neural} summarizes their key characteristics with respect to latency, naturalness, and real-time suitability.

\begin{table}[!htbp]
  \centering
  \caption[Traditional and neural TTS comparison]{Comparison between traditional and neural text-to-speech approaches in live interpretation scenarios}
  \label{tab:tts_traditional_vs_neural}
  \begin{adjustbox}{max width=\textwidth, keepaspectratio=false}
    \begin{tabular}{p{4.3cm} p{7.5cm} p{7.5cm}}
      \toprule
      \textbf{Aspect} & \textbf{Traditional TTS} & \textbf{Neural TTS} \\
      \midrule
      Synthesis approach
      &
      \begin{itemize}[nosep,leftmargin=*]
        \item Concatenative or parametric methods
        \item Predefined acoustic units
        \item Limited modeling of prosody
      \end{itemize}
      &
      \begin{itemize}[nosep,leftmargin=*]
        \item Data-driven neural models
        \item End-to-end acoustic modeling
        \item Improved prosody and expressiveness
      \end{itemize} \\
      \midrule
      Speech naturalness
      &
      \begin{itemize}[nosep,leftmargin=*]
        \item Robotic or monotonous output
        \item Limited speaker variation
      \end{itemize}
      &
      \begin{itemize}[nosep,leftmargin=*]
        \item Human-like speech quality
        \item Multiple voices and speaking styles
      \end{itemize} \\
      \midrule
      Latency characteristics
      &
      \begin{itemize}[nosep,leftmargin=*]
        \item Predictable and low synthesis delay
        \item Suitable for strict real-time constraints
      \end{itemize}
      &
      \begin{itemize}[nosep,leftmargin=*]
        \item Higher computational cost
        \item Requires latency-aware configuration
      \end{itemize} \\
      \midrule
      Suitability for live interpretation
      &
      \begin{itemize}[nosep,leftmargin=*]
        \item High responsiveness
        \item Limited listener comfort
      \end{itemize}
      &
      \begin{itemize}[nosep,leftmargin=*]
        \item Improved intelligibility and comfort
        \item Requires careful latency management
      \end{itemize} \\
      \bottomrule
    \end{tabular}
  \end{adjustbox}
\end{table}


The use of neural TTS in live interpretation introduces an inherent trade-off between speech naturalness and synthesis latency. Generating highly expressive speech typically requires additional processing time, which may conflict with the immediacy required for simultaneous interpretation. Consequently, real-time systems often prioritize fast synthesis and stable audio delivery over maximum expressiveness. This balance ensures that translated speech remains synchronized with the source speaker, preserving the simultaneity essential for user comprehension.

In addition to latency considerations, audio format compatibility plays a critical role in live interpretation scenarios. TTS systems must produce audio in standardized formats that can be streamed, buffered, and played back with minimal processing overhead. Raw pulse-code modulation (PCM) audio is commonly favored in real-time systems due to its simplicity and compatibility with professional audio pipelines. Compressed formats, while bandwidth-efficient, may introduce decoding delays that are undesirable in low-latency applications.

From a system-level perspective, TTS in live interpretation must integrate seamlessly with downstream audio routing and playback components. This includes supporting consistent sampling rates, channel configurations, and timing behavior across successive synthesized segments. Any variability in output timing or format can accumulate over time, leading to perceptible delays or playback artifacts as mentioned in Chapter~\ref{chapter:hardware}.



\section{AI-Based Simultaneous Interpretation Systems}

Recent advances in speech and language processing have enabled the development of fully AI-based simultaneous interpretation systems that aim to replicate, and in some contexts augment, the capabilities of traditional human-centered interpretation workflows. Unlike conventional systems that rely on human interpreters and specialized hardware, AI-based approaches attempt to automate the interpretation process by combining speech recognition, machine translation, and speech synthesis into an integrated computational pipeline \cite{popescu-belis2025s2streview}. These systems promise improved scalability, reduced operational costs, and increased accessibility for multilingual communication \cite{ozkaya2023transformative}.

Early research in automated interpretation explored tightly coupled, end-to-end models that directly map source-language speech to target-language speech. Such approaches seek to minimize latency by avoiding intermediate textual representations and reducing pipeline complexity. However, fully end-to-end speech-to-speech translation models remain challenging to deploy in real-time settings due to their high computational demands, limited controllability, and difficulty in handling incremental input \cite{jia2022translatotron2}. As a result, many practical systems continue to adopt cascaded architectures composed of separate automatic speech recognition (ASR), neural machine translation (NMT), and text-to-speech (TTS) components \cite{popescu-belis2025s2streview, Sethiya2023EndtoEndSurvey}. To summarize the architectural differences between end-to-end and cascaded AI-based interpretation systems, Table~\ref{tab:ai_interpretation_architectures} compares their key characteristics with respect to modularity, latency control, and practical deployment.

\begin{table}[!htbp]
  \centering
  \caption[AI-based interpretation architectures]{Comparison of end-to-end and cascaded AI-based simultaneous interpretation architectures}
  \label{tab:ai_interpretation_architectures}
  \begin{adjustbox}{max width=\textwidth}
    \begin{tabular}{p{4cm} p{6.5cm} p{6.5cm}}
      \toprule
      \textbf{Aspect} & \textbf{End-to-End Speech-to-Speech} & \textbf{Cascaded ASR--NMT--TTS} \\
      \midrule
      Architectural structure &
      Single unified model &
      Modular pipeline with separate components \\
      \midrule
      Latency control &
      Limited explicit control &
      Fine-grained, stage-level control \\
      \midrule
      Modularity &
      Low &
      High (components interchangeable) \\
      \midrule
      Error handling &
      Implicit and opaque &
      Explicit and stage-specific \\
      \midrule
      Practical deployment &
      Experimental &
      Widely adopted in practice \\
      \bottomrule
    \end{tabular}
  \end{adjustbox}
\end{table}

Cascaded interpretation pipelines offer several practical advantages over end-to-end models. By decoupling recognition, translation, and synthesis, each component can be optimized independently and replaced as improved models or services become available. This modularity also enables better error handling and system-level control, particularly in real-time scenarios where partial outputs and incremental updates must be managed carefully. Consequently, cascaded architectures remain the dominant paradigm in both academic prototypes and industrial systems for AI-based simultaneous interpretation \cite{popescu-belis2025s2streview}.

Academic research has extensively investigated the challenges associated with real-time interpretation, particularly the need to translate speech incrementally as it is spoken. A central focus of this work has been the design of strategies that determine when sufficient source-language information is available to trigger translation. Approaches such as wait-based policies, adaptive latency control, and incremental decoding have been proposed to balance translation quality against responsiveness. These methods typically operate on textual representations derived from streaming ASR output and aim to minimize delay while preserving semantic coherence \cite{ma2019staclarxiv, niehues2018simultaneousarxiv}.

Several research prototypes have demonstrated the feasibility of simultaneous translation under controlled conditions \cite{cho2016anticipation, weller2021streamingST, popescu-belis2025s2streview}. These systems often introduce intermediate buffering or anticipation mechanisms that allow translation to begin before full sentence completion. While such approaches reduce latency, they also introduce new sources of instability, including partial hypothesis revision and semantic inconsistency. As a result, system-level coordination between recognition, translation, and synthesis stages remains a critical challenge.

In parallel with academic research, industrial efforts have led to the emergence of cloud-based interpretation platforms that offer real-time translation services via web or application interfaces. These platforms typically rely on scalable cloud infrastructure and expose interpretation functionality through streaming APIs. While commercial systems demonstrate impressive performance and broad language coverage, they often operate as closed ecosystems, limiting transparency and flexibility. Moreover, their internal handling of latency, segmentation, and stability is rarely documented in detail, making systematic evaluation difficult.

A common characteristic across both academic and industrial systems is the reliance on networked, distributed architectures. Audio input is captured on a client device, transmitted to backend services for processing, and returned as translated text or synthesized speech. This distributed setup introduces additional latency due to network communication and requires persistent, low-latency connections to maintain continuous data flow \cite{he2019streamingarxiv}. Consequently, real-time interpretation systems must address not only model-level performance but also system-level concerns such as concurrency, buffering, and fault tolerance.

Typical AI-based interpretation architectures follow a streaming pipeline structure in which audio is processed incrementally and passed through successive transformation stages. Each stage produces intermediate representations that may be provisional or finalized, depending on available context. The timing and stability of these representations play a central role in determining overall system behavior \cite{ma2019staclarxiv}. In particular, the interaction between streaming ASR output and translation triggering mechanisms has been identified as a key factor influencing both latency and output quality.

Despite significant progress, existing AI-based interpretation systems continue to face limitations that hinder their deployment in real-world settings. Many systems struggle to maintain stable output under long-duration sessions, where latency accumulation and buffering effects can degrade performance over time. Others lack mechanisms for integrating with professional domain hardware, restricting their use to consumer-grade environments. Furthermore, provider-locked architectures reduce flexibility and complicate comparative evaluation across different AI services.



\section{Limitations of Existing Approaches and Research Gap}

The review of traditional and AI-based simultaneous interpretation systems highlights substantial progress toward automated multilingual communication\cite{ozkaya2023transformative, papi-etal-2022-laal}. Nevertheless, existing approaches continue to exhibit limitations that restrict their applicability in real-world, real-time settings. These limitations arise not only from model-level constraints but also from system-level design choices, particularly in the handling of streaming data, latency accumulation, and integration with practical deployment environments.

The fundamental challenge across AI-based interpretation systems is the accumulation of end-to-end latency \cite{zheng-etal-2020-fluent, ma2019staclarxiv, dhawan2022speech2speechchallenges}. While individual components such as speech recognition, translation, and speech synthesis may operate efficiently in isolation, their sequential combination often results in perceptible delays. In real-time interpretation scenarios, even small delays introduced at each stage can compound, leading to output that lags behind the source speech and disrupts the simultaneity expected by users.

Closely related to latency accumulation is the problem of segmentation which builds the basis of the translation quality. Many existing systems rely on simplistic segmentation strategies and insufficient stability check algorithms, which are poorly suited to streaming settings where streaming ASR systems continuously revise their hypotheses as new audio becomes available \cite{cho2016anticipation, chen2024metamorphosisMT}. These approaches can result in either premature segmentation that harms semantic coherence or delayed segmentation that increases latency. The lack of robust, stability-aware segmentation mechanisms remains a significant limitation in current systems.

Beyond algorithmic concerns, hardware integration presents an additional challenge \cite{conference_interpreting_new_tech}. While traditional interpretation systems are designed around professional audio infrastructure, many AI-based solutions are confined to software-only or consumer-oriented environments. Limited consideration is given to audio routing, standardized formats, or compatibility with interpretation consoles and transmission systems \cite{conference_interpreting_new_tech}. This gap hinders the adoption of AI-based interpretation in professional contexts where hardware integration is essential.

Finally, a notable limitation of many existing interpretation systems lies in their restricted flexibility for multi-language support \cite{weller2021streamingST, popescu-belis2025s2spipelines }. In practice, language coverage is often achieved by tightly coupling language-specific ASR, NMT, and TTS components to a single provider or fixed service configuration. As a result, many commercial platforms operate as closed, provider-locked ecosystems, limiting the ability to selectively combine or replace individual components to achieve optimal performance for different language pairs. This lack of modularity complicates adaptation to heterogeneous deployment scenarios and prevents systematic evaluation and comparison of alternative AI services across multiple languages.

% To clarify how the identified limitations are addressed within the scope of this thesis, Table~\ref{tab:limitations_mapping} summarizes the relationship between existing challenges and the corresponding design aspects considered.

% \begin{table}[!htbp]
%   \centering
%   \caption[Limitations of existing systems and addressed aspects]{Identified limitations of existing AI-based interpretation systems and corresponding aspects addressed in this thesis}
%   \label{tab:limitations_mapping}
%   \begin{adjustbox}{max width=\textwidth}
%     \begin{tabular}{p{6cm} p{11.3cm}}
%       \toprule
%       \textbf{Identified Limitation} & \textbf{Addressed Aspect in This Thesis} \\
%       \midrule
%       End-to-end latency accumulation &
%       Latency-aware pipeline design with asynchronous processing \\
%       \midrule
%       Poor or rigid segmentation strategies &
%       Stability-aware, clause-level segmentation mechanism \\
%       \midrule
%       Instability of partial ASR hypotheses &
%       Explicit stability filtering before translation triggering \\
%       \midrule
%       Limited hardware integration &
%       Hardware compatibility and feasibility considerations \\
%       \midrule
%       Provider-locked system architectures &
%       Modular, provider-agnostic component abstraction \\
%       \midrule
%       Lack of system-level evaluation &
%       Comprehensive end-to-end and component-level evaluation framework \\
%       \bottomrule
%     \end{tabular}
%   \end{adjustbox}
% \end{table}


Taken together, these limitations reveal a clear research gap at the intersection of real-time processing, system architecture, and practical deployment. Addressing this gap requires an interpretation framework that explicitly manages latency, stabilizes streaming outputs, supports modular component integration, and considers hardware compatibility from the outset. In Chapter~\ref{chapter:methodology} a methodology and system design that aims to address these challenges through a unified, real-time interpretation pipeline is proposed.

